\documentclass[11pt,a4paper]{article}
\usepackage[utf8]{inputenc}
\usepackage[T1]{fontenc}
\usepackage[margin=2.5cm]{geometry}
\usepackage{array}
\usepackage{booktabs}

\newcolumntype{L}[1]{>{\raggedright\arraybackslash}p{#1}}

\setlength{\parindent}{0pt}
\setlength{\parskip}{6pt plus 1pt minus 1pt}

\begin{document}

{\Large\textbf{Predicting Tool Confusion from Representational Separability in LLM Agents}}

\section*{1. Research question}

\textbf{To what extent does the separation between two tool specifications in a language model's
internal representation predict which of those tools the model will confuse when selecting?}

The four experiments form one chain. First establish what the model actually confuses, then check
whether that is visible internally, then find what causes it, and finally test whether the
diagnosis helps in practice.

\section*{2. Background}

An LLM agent picks a tool by reading a list of specifications. Each specification consists of a
name, a natural language description, and a parameter schema. When two specifications are similar,
the model sometimes selects the wrong one. Nothing in the pipeline flags this, since the call
still executes and returns data, so the error only surfaces downstream, if at all.

A documented pair from Norkart's Adresse- og eiendomss{\o}k API, taken from its public developer
documentation.

\begin{verbatim}
suggest_kommunecustom(query)
  "Immediate responses based on the entered search text.
   Not designed to handle typos."

search_kommunecustom(query)
  "More detailed information, with support for fuzzy searches
   that allow simple typos."
\end{verbatim}

The two endpoint names differ by a single token, the description vocabulary overlaps almost
fully, and both accept the same single text argument. Their behaviour differs in a way that
matters. The strict variant stops matching when the query contains a typo, while the fuzzy
variant absorbs it, so picking the wrong one changes results without raising an error.

\subsection*{2.1 Established premises}

Four published results are treated as premises rather than contributions.

\begin{enumerate}
    \item \textbf{Surface metadata causally drives selection.} Adversarial optimisation of a tool's name and
    description raises its selection rate from roughly 20\% to as high as 81\% (ToolTweak,
    2510.02554).
    \item \textbf{Selection bias is consistent across models.} Metadata-driven tool-selection bias is
    reproducible across frontier models (BiasBusters, 2510.00307).
    \item \textbf{Tool-related decisions are decodable from hidden states.} A linear probe on the
    last-input-token hidden state predicts whether a tool call is necessary with AUROC 0.89 to
    0.96, exceeding the model's own verbalised reasoning (When2Tool, 2605.09252).
    \item \textbf{Parameter schemas drive invocation independently of semantics.} Models invoke a tool
    whenever query attributes can be validly assigned to its parameters, even when the tool does
    not serve the goal. Contrastive attention attribution reveals competing semantic-checking and
    structural-matching pathways (Structural Alignment Bias, 2604.11322, ACL 2026).
\end{enumerate}

\subsection*{2.2 The gap}

A 2026 line of work has established that tool-related information is readable from hidden states.
Wu et al.\ (2605.07990) showed that the choice of tool is carried by a linear direction, one
direction per pair of tools, that probes trained on pre-generation states read which tool the
model will call at 61 to 89\% top-1 accuracy within a single domain, and that queries where the
model is unsure between two tools fail 21 times more often. Healy et al.\ (2601.05214) detect
hallucinated tool calls from hidden states in the same forward pass, and Chen (2606.16364)
localises selection failures to the decision readout rather than to attention over the tool list.

None of this work measures pairwise confusability. Wu et al.\ measure uncertainty per query and
pool across queries. They do not compute a separability score for a specific pair of tools, and
they do not test whether that score predicts the pair's confusion rate beyond what string
similarity or sentence embeddings already predict. Their unit of analysis is the query. Ours is
the tool pair. Healy et al.\ classify calls as correct or hallucinated without decoding tool
identity. Chen predicts per-query failure, not per-pair confusion. Premises (3) and (4) are both
about invocation, meaning whether to call a tool at all. Premise (1) treats the specification as
an optimisation target rather than a measurable object. Premise (2) concerns preference, which is
a different quantity from pairwise confusability. A tool can be systematically over-selected
without being confusable with any particular sibling.

What is missing is a measurement of whether internal geometry predicts pairwise selection errors,
and that measurement is the main goal of this thesis.

\subsection*{2.3 Closest related work}

\begin{tabular}{@{}L{3.2cm}L{5.9cm}L{6.2cm}@{}}
\toprule
\textbf{Work} & \textbf{What it does} & \textbf{Delta} \\
\midrule
Wu et al.\ (2605.07990) & Per-pair linear directions switch which tool the model picks, probes read which tool will be called, unsure queries fail 21x more often & Closest prior work. Anticipates RQ1. Uncertainty is measured per query and pooled, never per pair, and no surface-similarity baselines are tested. RQ2 is the pair-level, baseline-controlled version of their aggregate finding \\
When2Tool (2605.09252) & Probes hidden states for \textbf{whether} a tool call is needed & Binary necessity, not \textbf{which} tool, and no pairwise measure \\
Healy et al.\ (2601.05214) & Detects hallucinated tool calls from hidden states in one forward pass & Binary correct-versus-hallucinated classification, no identity decoding and no pairs \\
Chen (2606.16364) & Shows models attend to the correct tool and still pick wrong, and predicts per-query failure from attention margins against lexical and probability baselines & Per-query failure prediction, not per-pair confusion, and no pairwise identity separability \\
Tool-call dependency probing (2605.25310) & Linear probes decode dependency edges between sequential tool calls from the residual stream & Follows a call sequence over time, not identity separation between competing candidates, and no confusion prediction \\
Structural Alignment Bias (2604.11322) & Mechanistic analysis of invoke-versus-refuse, with attention attribution & Invocation, not discrimination between candidates. Motivates the schema factor in Study 3 \\
ToolTweak (2510.02554) & Adversarially optimises metadata to raise selection share & No representational measurement and no confusion prediction \\
BiasBusters (2510.00307) & Tool preference bias and cross-model consistency, with metadata perturbations & Preference bias is not confusability, and the perturbations are one factor at a time with no interaction terms \\
Canary Tools (2608.04719) & Six-type taxonomy of planted diagnostic tools that probe specific selection weaknesses & Behavioural decoy taxonomy, no internal measurement \\
ToolSandbox (2408.04682) & Benchmark with name scrambling, description scrambling, and similarity-selected distractors & Scrambles one component at a time, reports task accuracy rather than pairwise confusion, and has no factorial crossing or interaction analysis \\
\bottomrule
\end{tabular}

Experiment 3 is a partial replication and extension rather than a novel design. It is kept because it
makes Experiment 2 measurement actionable, and because published evidence on the question
conflicts. The novelty rests on RQ2.

\subsection*{2.4 Why open weights are the right setting}

The method needs access to model internals, which is exactly what self-hosting provides. Capable
open-weight dense models now exist in a size range that runs comfortably in bf16 on modest
hardware. For an organisation running its own model, the detector probes the same weights,
tokeniser, and specifications that serve production, so no transfer assumption is involved.

\section*{3. Research questions}

\textbf{RQ1.} To what extent, and at what network depth, is the identity of the correct tool linearly
decodable from a model's pre-generation hidden states?

\textbf{RQ2 (central).} To what extent does pairwise representational separability between two tools
predict the rate at which the model confuses them, relative to what surface-similarity and
sentence-embedding baselines already predict?

\textbf{RQ3.} To what extent is representational separability attributable to each component of a
specification, meaning name, description, and parameter schema, and under what conditions does
each dominate?

The published evidence points in different directions. On the side of names, ToolTweak found that
two identical tools differing only by a numeric suffix received very different selection rates,
and PA-Tool (2510.07248) reports gains of up to 17\% from renaming tools to pretraining-aligned
names. On the side of descriptions, BiasBusters found name-only perturbations small and
high-variance while description semantics were the primary discrimination cue, and ToolSandbox's
GPT-4o ablations dropped 3.7 points under description scrambling against 0.6 under name
scrambling. We therefore expect no main effect of any single component. Instead we expect an
interaction between name and description, where name overlap matters most when descriptions leave
the boundary between two sibling tools unstated, and description quality matters most when the
names are already distinct.

\textbf{RQ4.} To what extent do specification edits guided by the RQ3 attribution reduce measured
confusion on a real tool set, relative to a sham edit of equal effort on the same pair?

\section*{4. Method}

\subsection*{Experiment 1. Behavioural ground truth}

We build a query set over the Norkart tool inventory in which every query has one unambiguous
correct tool. Two independent annotators verify the labels. Inter-annotator agreement is
reported, ambiguous queries are excluded, and the excluded count is reported alongside.

The primary dependent variable is the probability mass assigned to each competing tool, taken
from the tool-choice distribution on the same forward pass. Argmax selection accuracy is kept as
a secondary metric because it is easy to interpret. Probability mass is the better construct.
Argmax makes confusion a rare binary event and leaves most pair-level cells near zero, while
probability mass is continuous and far better powered. A pair where the model puts 45\% on the
wrong tool is confusable even when it happens to select correctly.

Decoding settings are pinned before data collection. Greedy decoding, temperature 0, a fixed
repeat count per query, and unconstrained JSON output parsed afterwards. Tool-list order is
randomised per trial and reported as a covariate.

\subsection*{Experiment 2. Representational separability and prediction (RQ1, RQ2)}

Residual-stream activations are extracted at the final prompt token, across all layers, in bf16,
with thinking disabled.

\begin{enumerate}
    \item \textbf{Decodability (RQ1).} Multiclass linear probe predicting the correct tool, swept across layers.
    This step replicates and extends Wu et al.\ (2605.07990). The extension is the label. Their
    probes predict which tool the model will call, while ours predict which tool is the correct
    target, and the gap between the two is itself diagnostic of confusion.
    \item \textbf{Pairwise separability (RQ2).} For each pair, a two-class probe separates queries targeting
    one tool from queries targeting the other. The primary separability measure is the
    cross-validated accuracy of that probe. Centroid cosine distance and a between/within-class
    scatter ratio are reported alongside as descriptive support.
    \item \textbf{Prediction.} Correlate separability against Experiment 1 confusion measure and test
    incremental variance over the baselines in \S5.
\end{enumerate}

Three constraints address specific ways this measurement can go wrong.

\begin{itemize}
    \item \textbf{Template-matched queries.} A two-class probe trained on two different query subsets can
    separate on stylistic register rather than tool identity. All Experiment 2 queries use a shared
    frame with a varying slot fill, so the only systematic difference between subsets is the target
    tool.
    \item \textbf{Position-balanced extraction.} Activations are position-sensitive, so extraction runs are
    balanced across tool-list positions. Without this, separability partly encodes list position
    rather than tool identity.
    \item \textbf{Raw and residualised separability both reported.} Residualising query length, token count,
    and lexical overlap with the tool name is necessary because probe accuracy alone is not
    evidence of internal structure. Recent work found probe separations collapsing to chance once
    format confounds were removed (arXiv 2606.02907). Lexical overlap is also the mechanism RQ3
    hypothesises, so residualising it away answers a different question. Raw separability answers
    whether geometry is predictive. Residualised separability answers whether geometry adds
    anything beyond strings. Both are reported, with what each licenses stated explicitly. A
    permutation baseline with shuffled labels accompanies every probe result, and the primary
    measurement is taken mid-stack with the full layer sweep reported.
\end{itemize}

\subsection*{Experiment 3. Component attribution (RQ3)}

A factorial design over the three specification components, run on matched Norkart-domain pairs.

\begin{tabular}{@{}L{2.8cm}L{5.9cm}L{6.5cm}@{}}
\toprule
\textbf{Factor} & \textbf{Level A} & \textbf{Level B} \\
\midrule
Name & shared stem (\texttt{eiendomsflater} / \texttt{eiendomskart}) & lexically disjoint (\texttt{eiendomsflater} / \texttt{ruteberegner}) \\
Description & independent paraphrase & contrastive (states the boundary against the sibling) \\
Parameter schema & identical arity and names & divergent arity or types \\
\bottomrule
\end{tabular}

Dependent variables are the confusion measure and representational separability. Effect sizes
come with confidence intervals. The name $\times$ description interaction is the pre-registered focus.
The schema factor is motivated directly by the structural-matching pathway identified in premise
(4).

Name similarity is measured in each model's own token vocabulary rather than in characters, since
the hypothesised mechanism is shared token prefixes.

\textbf{Norwegian.} Norwegian compounds fragment differently under model tokenisers, which supplies
natural variation in the name factor. This arm runs on Qwen3-14B only, because Granite 4.1 does
not list Norwegian among its supported languages. Granite contributes to the English cells and to
RQ1 and RQ2, not to the bilingual comparison.

\subsection*{Experiment 4. Intervention (RQ4)}

This experiment runs on Norkart's real tool set. The Experiment 2 detector flags risky pairs. For each
flagged pair we produce two edits of matched length and effort. The guided edit targets whatever
cause Experiment 3 attributes. The sham edit is an equally careful rewrite that does not target that
cause. Every flagged pair is then measured under all three specifications, original, guided, and
sham, in separate balanced trials with the condition order randomised.

The rule that maps an Experiment 3 attribution to a concrete edit is fixed before Experiment 4 begins. A
shared name stem leads to renaming toward a lexically distinct stem. A weak boundary leads to a
contrastive description that states the boundary. Schema-driven confusion leads to divergent arity
or types.

Because every flagged pair serves as its own control, the design avoids the
regression-to-the-mean objection that comparing flagged against unflagged pairs would invite.

\section*{5. Evaluation}

Three baselines are committed in advance.

\begin{enumerate}
    \item \textbf{String similarity.} Normalised edit distance and token-set overlap on names, computed in
    the model's tokeniser vocabulary.
    \item \textbf{Sentence embedding.} Cosine similarity of full specifications under a multilingual encoder.
    \item \textbf{Output-side pair margin.} For each pair we take the queries in which one of the two tools
    is the correct target and record how much probability mass lands on each of the two. A small
    gap between the two masses marks the pair as confusable. This signal comes from the
    tool-choice distribution alone and needs no internal access.
\end{enumerate}

The primary RQ2 metric is incremental variance explained over baselines 1 to 3, plus precision
and recall of the top-k flagged pairs against the pairs that actually produce errors.

\textbf{Sample size and power.} Queries per tool, trials per query, and a minimum-observation
inclusion rule per pair are fixed before data collection. A power statement for the RQ3
interaction is included, since detecting an interaction requires roughly four times the sample of
a main effect of equal size, and this drives the design. The continuous confusion measure from
Experiment 1 substantially reduces the required sample compared to a binary argmax outcome.

\textbf{Statistical treatment.} Tool pairs are not independent because each tool appears in many
pairs. Analyses therefore use mixed-effects models with crossed random intercepts per tool, or
cluster-robust standard errors. Confusion proportions are modelled with beta-binomial or
equivalent overdispersion handling.

\textbf{Negative-result plan.} If the internal measures do not improve on baselines 1 to 3, the answer
to RQ2 is that surface string overlap suffices and internals are unnecessary for this task.
Experiments 3 and 4 still stand under that outcome, with the detector simply running on the cheaper
signal. The thesis therefore produces a usable result under either answer to its central question.

\section*{6. Models and materials}

\subsection*{6.1 Tool inventory}

Norkart's published services supply non-synthetic semantic overlap. Adresse- og eiendomss{\o}k
(autocomplete, fuzzy, reverse geocoding), Eiendomsflater, Naboliste, Eiendom og Tinglysning,
Bygning, Takgeometri og solinnstr\r{a}ling, Eiendomskart, Geografisk h{\o}yde, Ruteberegner,
administrative units, and Bakgrunnskart.

Confusable pairs available without construction.

\begin{itemize}
    \item \texttt{suggest\_kommunecustom} vs.\ \texttt{search\_kommunecustom}. Two documented endpoints of one service,
    same single string input, overlapping description vocabulary, strict versus fuzzy matching
    \item reverse geocoding vs.\ the text based modes of the same service. Overlapping intent, coordinate
    input instead of text
    \item \texttt{matrikkelenhet} / \texttt{grunneiendom} / \texttt{festegrunn} / \texttt{eierseksjon} / \texttt{anleggseiendom}. Legally
    distinct cadastral objects with overlapping lexical form
    \item \texttt{eiendomsflater} vs.\ \texttt{eiendomskart}. Two separate documented services, a one-token difference,
    parcel attribute data versus property boundary map service
\end{itemize}

The experiment needs specifications, not data. No live API calls and no cadastral records are
required. The object of study is which tool the model selects, not what it returns.

\subsection*{6.2 Models}

Models are chosen for measurement fidelity rather than maximum capability. They must be small
enough to run in bf16, dense and standard enough for clean per-layer hidden-state extraction, and
drawn from two families so the findings are not tied to one architecture.

\begin{tabular}{@{}L{1.9cm}L{4.3cm}L{9.2cm}@{}}
\toprule
\textbf{Role} & \textbf{Model} & \textbf{Rationale} \\
\midrule
Primary A & Granite 4.1 8B (dense, Apache 2.0) & Strong function calling for its size, and no reasoning mode, matching \S6.3 \\
Primary B & Qwen3-14B (dense, Apache 2.0) & Standard full-attention transformer for clean per-layer extraction, supports 119 languages including Norwegian Bokm{\aa}l and Nynorsk, carries the Norwegian arm of Study 3 \\
Contingent & Qwen3-32B (dense, Apache 2.0) & Deployment-scale spot check, same dense architecture family as Primary B. Runs in bf16 on a single 80\,GB card at moderate context with hidden states offloaded per batch \\
Appendix & One closed model with logprob access & Behavioural robustness only, no internal claim \\
\bottomrule
\end{tabular}

All activation extraction runs in bf16. Quantised activations are never used for measurement,
because quantisation noise falls directly on the small inter-tool distances that are the
dependent variable.

\subsection*{6.3 Inference mode}

All internal measurement runs with thinking disabled. For Qwen3 this is set through the
\texttt{enable\_thinking=False} flag in the chat template, since the model reasons by default.
Granite 4.1 has no reasoning mode, so no setting is needed there. The probe assumes tool identity
is resolved at the last prompt token. If a model reasons for thousands of tokens before
selecting, the decision is made inside the trace and the probe location is invalid. Extending
this analysis to reasoning models would require locating the decision point within the trace,
which we state as future work.

\section*{7. Contributions}

\begin{enumerate}
    \item \textbf{Primary.} To my knowledge, as of August 2026, this is the first work to measure
    representational separability per tool pair and to test whether it predicts that pair's
    confusion rate beyond surface-similarity, sentence-embedding, and output-margin baselines.
    The closest prior work, Wu et al.\ (2605.07990), established that tool choice is linearly
    readable and that uncertain queries fail more often, but measured uncertainty per query and
    pooled across queries, with no per-pair separability statistic and no surface baselines. RQ2
    is the pair-level, baseline-controlled analysis their aggregate finding suggests but does not
    perform. RQ1 replicates their decodability result and extends it from will-call to
    correct-target decoding.
    \item Resolution of conflicting published evidence on whether names or descriptions dominate
    confusability, framed as an interaction rather than a main effect.
    \item Causal evidence, against a within-pair sham control, that attribution-guided edits reduce
    selection errors.
    \item An open-source design-time confusability detector returning ranked risk pairs with attributed
    cause.
    \item A public tool set and query set in the Norwegian cadastral domain.
\end{enumerate}

\section*{8. Principal risks}

\begin{tabular}{@{}L{5.4cm}L{10.0cm}@{}}
\toprule
\textbf{Risk} & \textbf{Mitigation} \\
\midrule
Internal measures do not beat string baselines & Pre-committed negative-result framing (\S5), Experiments 3 and 4 stand alone \\
Probe signal is a confound artefact & Template-matched queries, position balancing, residualisation, permutation baselines \\
Sparse confusion matrix undermines correlations & Continuous probability-mass DV, inclusion rule and power statement fixed in advance \\
A competing preprint appears mid-project & arXiv alerts from week 1, Experiment 4 and the Norkart tool set are not scoopable by benchmark-only work \\
Qwen3-32B does not fit the 80\,GB card in bf16 & Contingent model only, falls back to a behavioural spot check at no cost to the chain \\
Norkart access delayed & Experiments 1 to 3 run on specifications authored from public documentation \\
\bottomrule
\end{tabular}

\section*{9. References}

\begin{itemize}
    \item Alain \& Bengio (2016). Understanding intermediate layers using linear classifier probes.
    \item Belinkov (2022). Probing classifiers: promises, shortcomings, and advances.
    \item Patil et al. (2023). Gorilla: large language model connected with massive APIs.
    \item Qin et al. (2023). ToolLLM.
    \item Yan et al. (2024). Berkeley Function Calling Leaderboard.
    \item ToolSandbox (arXiv 2408.04682).
    \item ToolTweak (arXiv 2510.02554).
    \item BiasBusters (arXiv 2510.00307).
    \item PA-Tool (arXiv 2510.07248).
    \item Do LLMs know tool irrelevance? Structural alignment bias in tool invocations (arXiv 2604.11322).
    \item LLM agents already know when to call tools (arXiv 2605.09252).
    \item Tool-call dependency structure is linearly decodable in LLM agent residual streams (arXiv 2605.25310).
    \item Tool calling is linearly readable and steerable in language models (arXiv 2605.07990).
    \item Internal representations as indicators of hallucinations in agent tool selection (arXiv 2601.05214).
    \item Looking is not picking: an attention-segment account of tool-selection failures in agents.
    arXiv 2606.16364.
    \item Linear probes detect task format, not reasoning mode (arXiv 2606.02907).
    \item Diagnosing tool-selection reasoning with canary tools (arXiv 2608.04719).
    \item IBM Research (2026). Granite 4.1 family of models.
    \item Norkart AS. Datatjenester API documentation.
\end{itemize}

\end{document}